% Compile no Overleaf com pdfLaTeX.
% Preencha os campos da capa antes da entrega.
\documentclass[10pt,a4paper]{article}

\usepackage[utf8]{inputenc}
\usepackage[T1]{fontenc}
\usepackage[brazil]{babel}
\usepackage{lmodern}
\usepackage{geometry}
\usepackage{array}
\usepackage{tabularx}
\usepackage{booktabs}
\usepackage[hidelinks]{hyperref}

\geometry{left=2.2cm,right=2.2cm,top=2.0cm,bottom=2.0cm}
\setlength{\parindent}{0pt}
\setlength{\parskip}{5pt}
\renewcommand{\arraystretch}{1.12}

\begin{document}

% --- Página 1: capa ---
\begin{titlepage}
\centering
{\large \textbf{Pontifícia Universidade Católica do Rio Grande do Sul}\par}
\vspace{0.5cm}
{\normalsize Engenharia de Software\par}
\vspace{0.5cm}
{\normalsize Disciplina: Coleta, Preparação e Análise de Dados\par}
\vfill
{\Large\bfseries Laboratório 1\par}
\vspace{0.35cm}
{\LARGE\bfseries Web Scraping de Vagas de Emprego\par}
\vspace{0.5cm}
{\large Portal ProgramaThor\par}
\vfill
\begin{minipage}{0.77\textwidth}
\textbf{Grupo:} Leonardo Fagundes Wingert, Gustavo de Freitas Fidelis e Samuel Bitdinger\par
\vspace{0.8cm}
\textbf{Professora:} Katherine Bianchini Esper de Vargas\par
\end{minipage}
\vfill
{Porto Alegre, RS\par}
{2026\par}
\end{titlepage}

% --- Página 2: desenvolvimento e resultados ---
\pagenumbering{arabic}
\section*{1. Objetivo e decisões de desenvolvimento}

Antes de escrever o scraper, avaliamos o Indeed. Sua análise preliminar de robots.txt mostrou restrições a /job e /viewjob, caminhos das páginas individuais necessárias para obter descrição e requisitos. Por isso, o portal foi descartado antes de qualquer implementação. Em seguida, iniciamos o código no Vagas.com.br: a busca funcionou, mas, durante os testes de acesso às vagas individuais, o site passou a mostrar a mensagem ``Sua solicitação foi bloqueada''. Esse limite foi encontrado na prática, durante o desenvolvimento, e motivou a troca de portal.

O objetivo final foi coletar vagas de desenvolvimento de software, na área de Python, em três páginas do \href{https://programathor.com.br/jobs}{ProgramaThor}. Ele oferece páginas individuais com atividades e requisitos e paginação convencional. A busca não começa em uma URL filtrada: o Selenium abre a lista geral e clica no link visível \texttt{Vagas programador PYTHON}, interagindo com o site como exige o enunciado.

\section*{2. Verificação de robots.txt e sitemap}

Antes da coleta, foi examinado o \href{https://programathor.com.br/robots.txt}{robots.txt}. Na consulta preparatória de 20/09/2026, o bloco \texttt{User-agent: *} restringia \texttt{/admin/}, \texttt{/user/}, \texttt{/users/} e \texttt{/company/}. As rotas necessárias --- \texttt{/jobs}, \texttt{/jobs-python} e \texttt{/jobs/...} --- não apareciam entre as restrições. A verificação desses arquivos foi feita fora do scraper, conforme a decisão do grupo.

O robots.txt indicava o \href{https://programathor.com.br/sitemap.xml}{sitemap.xml}. A cópia examinada continha 30.382 URLs, incluindo categorias por tecnologia, localidades e anúncios individuais. O sitemap serviu para conhecer a estrutura do portal; não foi usado como fonte de vagas ou como atalho para a busca.

\section*{3. Ferramentas, coleta e tratamento}

O Selenium automatiza o Safari: abre o portal, aciona o filtro, avança pelas páginas e visita cada anúncio. Esperas explícitas confirmam o conteúdo carregado, e há um intervalo de cinco segundos entre acessos. O BeautifulSoup interpreta o HTML das listas e das páginas individuais. A biblioteca padrão do Python normaliza URLs, trata datas e grava o \texttt{CSV}; nenhuma API de vagas foi utilizada.

Nas listas, o programa guarda título, link, página e selo de situação. O link identifica duplicatas antes de visitar as vagas. Nas páginas individuais, extrai título, empresa, local, salário, atividades/descrição e requisitos. Campos inexistentes recebem \texttt{Não informado}; expressões como \texttt{Não especificado} ou \texttt{Até R\$18.000} são mantidas como publicadas, sem inventar valores. Erros de navegação ou de extração são tratados; uma vaga com falha não interrompe as demais. O resultado é um único arquivo estruturado, \texttt{dados/vagas.csv}.

\section*{4. Resultado e limites}

O arquivo final contém \textbf{41 vagas} obtidas nas três páginas percorridas: 12 registros vieram da primeira página, 15 da segunda e 14 da terceira. A conferência do CSV não encontrou URLs repetidas nem campos vazios; título, empresa, local, salário, link, descrição e requisitos estão preenchidos em todas as linhas. Quando o portal não divulgava remuneração, o valor original \texttt{Não especificado} foi preservado. Isso ocorreu em 28 vagas; as outras 13 apresentavam alguma informação salarial.

Quanto ao local de trabalho, 22 vagas foram anunciadas como remotas e 19 como híbridas ou presenciais. O portal marcou 36 registros como \texttt{Vencida}; somente cinco não exibiam esse selo. A ausência do selo não comprova que a oportunidade ainda esteja aberta. Assim, o CSV atende à coleta estruturada e à navegação por três páginas, mas a grande quantidade de anúncios vencidos limita seu uso como relação de oportunidades atuais.

\clearpage

% --- Página 3: todos os métodos de busca e declaração de IA ---
\section*{5. Métodos de localização utilizados no código}

\small
\setlength{\tabcolsep}{4pt}
\begin{tabularx}{\textwidth}{@{}p{3.2cm}p{5.8cm}X@{}}
\toprule
\textbf{Método} & \textbf{Seletor / operação} & \textbf{O que localiza} \\
\midrule
Selenium \texttt{By.LINK\_TEXT} & \texttt{Vagas programador PYTHON} & Link de filtro por tecnologia na lista geral. \\
Selenium \texttt{By.TAG\_NAME} & \texttt{h1} & Cabeçalho que confirma a página de resultados Python. \\
Selenium CSS & \texttt{.pagination .active} & Número da página atual. \\
Selenium/BS4 CSS & \texttt{.cell-list > a[href]} & Links principais dos cards de vagas. \\
Selenium CSS & \texttt{.pagination a[rel="Próx"]} & Link para a página seguinte. \\
Selenium CSS & \texttt{link[rel="canonical"]} & URL canônica da vaga individual carregada. \\
Selenium CSS & \texttt{.wrapper-content-job-show .line-height-2-4} & Bloco que contém o texto da vaga. \\
BS4 \texttt{select\_one} & \texttt{h3}, dentro do card & Título mostrado na lista. \\
BS4 \texttt{select} & \texttt{span}, dentro de \texttt{h3} & Selos da listagem, inclusive \texttt{Vencida}. \\
BS4 \texttt{select\_one} & \texttt{.wrapper-content-job-show} & Contêiner dos detalhes da vaga. \\
BS4 \texttt{select\_one} & \texttt{.wrapper-header-job-show h1} & Título da página individual. \\
BS4 \texttt{select\_one} & \texttt{h2}, no contêiner de detalhes & Nome da empresa, com ou sem link. \\
BS4 \texttt{select} & \texttt{.wrapper-details-job-show p} & Parágrafos com local e salário, reconhecidos pelos rótulos \texttt{Localização:} e \texttt{Salário:}. \\
BS4 \texttt{select} & \texttt{.line-height-2-4 h3} & Títulos das seções de descrição e requisitos. \\
BS4 \texttt{find\_next\_siblings} & Irmãos até o próximo \texttt{h3} & Texto pertencente a cada seção da vaga. \\
\bottomrule
\end{tabularx}
\normalsize

\vspace{0.35cm}
\noindent\textbf{Nota:} CSS significa seletor CSS; BS4 significa BeautifulSoup. \texttt{get\_text()} e operações de strings limpam o texto depois da localização. O código não utiliza XPath nem expressões regulares para localizar elementos.

\section*{6. Observações sobre o uso de IA}

Foi utilizado o OpenAI Codex para auxiliar na escolha e análise do portal, no diagnóstico da automação do Safari, na identificação de seletores, na implementação e revisão do scraper, nos testes de extração, na conferência do CSV e na redação deste relatório. A IA também auxiliou a interpretar o enunciado e a registrar a análise preparatória de robots.txt e sitemap.

\end{document}
